\documentclass[11pt]{article}

\usepackage[margin=1in]{geometry}
\usepackage[T1]{fontenc}
\usepackage{listings}
\usepackage{xcolor}
\usepackage{amsmath}
\usepackage{amssymb}

\setlength{\parindent}{0pt}
\setlength{\parskip}{0.5em}

\lstset{
  language=C,
  basicstyle=\ttfamily\small,
  keywordstyle=\color{blue!60!black},
  commentstyle=\color{green!45!black}\itshape,
  showstringspaces=false,
  breaklines=true,
  frame=single,
  framesep=4pt,
  xleftmargin=4pt,
}

\begin{document}

\begin{center}
    {\bfseries
    Describe how you would add a copying garbage collector to the L3 virtual machine to replace the mark and sweep one that you implemented. In particular, explain whether the fact that a copying garbage collector has to rewrite pointers after copying objects would be problematic. If yes, explain the problem as well as potential solutions.
    }
\end{center}
\vspace{-2mm}
To replace the mark-and-sweep collector with a copying one, I would first reorganize the heap into two equal semi-spaces. At any point, the mutator allocates only in one of them (from-space), using a single \texttt{free} pointer bump. The other semi-space (to-space) is left empty until garbage collection (GC) is triggered. Collection copies every reachable block into to-space, swaps the roles of the two spaces, and resumes. Allocation is then as fast as stack allocation, GC cost is proportional to the live set, and the absence of a free list eliminates external fragmentation. The cost is that only half the heap is usable at a time and the GC must copy all the live data.

The \texttt{memory} struct would lose its bitmap, free lists, and sweep logic and gain four boundary pointers plus a bump pointer. Allocation then reduces to a bounds check and a pointer bump, with a collection triggered on overflow:

\begin{lstlisting}
value* memory_allocate(memory* self, tag tag, value size, value* root) {
  value total = HEADER_SIZE + size;
  if (self->free + total > self->from_end) {
    memory_collect(self, root);
    if (self->free + total > self->from_end)
      fail("no memory left (block of size %u requested)", size);
  }
  value* block = self->free + HEADER_SIZE;
  block_set_tag_size(block, tag, size);
  self->free += total;
  return block;
}
\end{lstlisting}
The GC starts from the roots. In the L3VM, the main root is the current top-frame block pointed to by \texttt{FP}. There is one special case: if the caller frame is stored in the other top-frame block, that second block must be scanned as well. These top-frame blocks live outside the heap, so the copying collector must not move them. It scans them in place and rewrites the heap pointers they contain. The copying pass itself can then use Cheney's algorithm. 

Pointer rewriting is \textbf{in general} problematic for a copying garbage collector. Before rewriting a word, the GC must know with certainty that the word is a pointer. If it rewrites a word that is actually an integer or a code address, that value is corrupted into a meaningless heap address. This is the key difference with conservative mark-and-sweep. In mark-and-sweep, mistaking a non-pointer for a pointer only causes the collector to mark some extra block as live. That is wasteful, but still safe. In a copying collector, the same mistake has bigger consequences, because the false pointer is followed and then overwritten with the new address of the moved block. 

At the L3 language level, values are tagged. Integers are encoded as \texttt{(n << 1) | 1}, i.e. as $2n+1$ modulo 32 bits, so their low bit is set. Characters, booleans, and unit also have non-zero low-bit patterns. Block pointers are word-aligned virtual addresses, so their two lowest bits are zero. This means that for ordinary L3 values, a cheap alignment test already rules out all immediates: \vspace{-0.5mm} \[ (v \mathbin{\&} (\texttt{VALUE\_BYTES} - 1)) \neq 0. \] \vspace{-0.5mm} Following the L3 value representation, such a value cannot be a block pointer and should not be evacuated. This test is only the first step, however. The L3VM is lower-level than the source language, and not every word stored in a frame is a tagged L3 value. In particular, untagged integers can appear in activation frames. Since activation frames may be evicted to the heap, the collector can encounter such words while scanning heap objects. An unlucky untagged integer may have the same bit pattern as a valid heap address. The bitmap used by the mark-and-sweep collector is enough to avoid treating a random address in the middle of the heap as an object, but it still does not prove that the original word was a pointer. It only proves that the word points to the beginning of some allocated block. 

There are also fixed raw slots in some heap objects. For example, slot~0 of a \texttt{Function} block is a code pointer, not an L3 heap pointer, and slot~0 of a \texttt{RegisterFrame} block is saved control information. These slots must not be passed blindly to \texttt{evacuate}, even if their low bits look pointer-like. 

So the GC should use two levels of precision. First, it uses the L3 tag bits to reject ordinary immediates. Second, it scans each block according to its layout. A string block contains no heap pointers. A function block skips its code pointer and scans only the captured values. Ordinary data blocks can be scanned field by field as L3 values. Register frames are the special case: because they may contain untagged temporaries, blindly scanning every register slot is not enough for a moving collector. Either the VM/compiler must provide a frame map saying which slots contain real heap pointers, or the implementation must enforce an invariant that all values live at a GC point are properly tagged L3 values.

\begin{lstlisting}
static void scan_block(memory* self, value* block, frame_map* map) {
  tag   t    = block_tag(block);
  value size = block_size(block);
  switch (t) {
  case tag_String:
    return;                       /* raw bytes, no pointers */

  case tag_Function:
    /* slot 0 is a raw code pointer; captured values start at slot 1 */
    for (value i = 1; i < size; ++i)
      block[i] = evacuate(self, block[i]);
    return;

  case tag_RegisterFrame:
    /* slot 0 is saved control information; slot 1 is the parent frame */
    if (size > 1)
      block[1] = evacuate(self, block[1]);

    /* Remaining slots must be scanned precisely, not conservatively. */
    for (value i = REGISTER_FRAME_CONTEXT_SIZE; i < size; ++i)
      if (frame_map_slot_is_pointer(map, i))
        block[i] = evacuate(self, block[i]);
    return;

  default:
    /* Ordinary L3 block: fields are tagged L3 values. */
    for (value i = 0; i < size; ++i)
      block[i] = evacuate(self, block[i]);
  }
}
\end{lstlisting}

With such a precise scan, pointer rewriting is safe: every word that is rewritten is known to be a real heap pointer, and every reachable heap pointer is eventually updated. Without that precision, copying is not safe. The alternatives are the standard ones: emit stack/frame maps and per-tag layout descriptors; introduce stable handles so program references do not change when objects move; or use a mostly-copying collector, where ambiguous objects are pinned and only unambiguous objects are moved. The important point is that the old conservative bitmap test from mark-and-sweep cannot simply be reused as the basis for moving objects. 

The full collection loop then follows Cheney directly: scan the pinned root frames in place, process to-space breadth-first until \texttt{scan} reaches \texttt{free}, and finally swap the two semi-spaces so the next allocation begins at the start of what used to be to-space. The loop itself is standard; the interesting part is making the scan precise enough for the pointer rewrites to be correct. 

In conclusion, the mechanics of the collector are clean: allocate linearly, copy the live graph to to-space, leave forwarding pointers behind, and flip the two semi-spaces. The real difficulty is not moving the blocks, but knowing exactly which words are pointers to moved blocks. Mark-and-sweep can survive a conservative guess, because a false pointer only keeps garbage alive. A copying collector cannot: a false pointer is rewritten, turning a harmless-looking ambiguity into data corruption. For the L3VM, a copying collector is therefore correct only if the scan is precise. Ordinary tagged L3 values and fixed block layouts get us most of the way there, but register frames require either precise frame maps, a strong ``all-live-values-are-tagged-at-GC'' invariant, or pinning of ambiguous objects. With that precision in place, copying GC is a good fit; without it, it is not safe.
\end{document}